\documentclass[11pt,a4paper]{article}
\usepackage[utf8]{inputenc}
\usepackage[T1]{fontenc}
\usepackage{amsmath,amssymb}
\usepackage{geometry}
\usepackage{enumitem}
\usepackage{multicol}
\usepackage{booktabs}

\geometry{margin=0.75in}

\title{\textbf{NEET Biology Practice Set: Meiosis I Stages}\\ \large Strictly Aligned with NCERT Class 11 (Chapter 10: Cell Cycle and Cell Division)}
\author{45 High-Yield MCQs with Answer Key \& Detailed Solutions}
\date{}

\begin{document}
\maketitle

\section*{Part I: Multiple Choice Questions (45 Questions)}

\begin{enumerate}[label=\textbf{\arabic*.}]

\item Identify the correct sequence of events occurring during Meiosis I strictly according to NCERT:
\begin{enumerate}[label=(\Alph*)]
    \item Synapsis $\rightarrow$ Crossing over $\rightarrow$ Terminalisation of chiasmata $\rightarrow$ Disjunction of homologs
    \item Crossing over $\rightarrow$ Synapsis $\rightarrow$ Terminalisation of chiasmata $\rightarrow$ Splitting of centromere
    \item Terminalisation of chiasmata $\rightarrow$ Synapsis $\rightarrow$ Crossing over $\rightarrow$ Disjunction of homologs
    \item Synapsis $\rightarrow$ Terminalisation of chiasmata $\rightarrow$ Crossing over $\rightarrow$ Disjunction of homologs
\end{enumerate}

\item Which of the following statements regarding Prophase I of Meiosis is \textbf{incorrect} as per NCERT?
\begin{enumerate}[label=(\Alph*)]
    \item It is typically longer and more complex when compared to prophase of mitosis.
    \item It has been subdivided into five phases based on chromosomal behaviour.
    \item During leptotene, chromosomes become gradually visible under the light microscope.
    \item Compaction of chromosomes stops immediately at the end of leptotene.
\end{enumerate}

\item Match Column I with Column II with respect to the sub-stages of Prophase I:
\begin{center}
\begin{tabular}{ll}
\toprule
\textbf{Column I} & \textbf{Column II} \\
\midrule
(a) Leptotene & (i) Action of enzyme recombinase \\
(b) Zygotene & (ii) Dissolution of synaptonemal complex \\
(c) Pachytene & (iii) Compaction of chromosomes continues throughout \\
(d) Diplotene & (iv) Formation of bivalent or tetrad \\
\bottomrule
\end{tabular}
\end{center}
\begin{enumerate}[label=(\Alph*)]
    \item (a)-(iii), (b)-(iv), (c)-(i), (d)-(ii)
    \item (a)-(iv), (b)-(iii), (c)-(i), (d)-(ii)
    \item (a)-(iii), (b)-(i), (c)-(iv), (d)-(ii)
    \item (a)-(ii), (b)-(iv), (c)-(i), (d)-(iii)
\end{enumerate}

\item Consider the following two statements:
\textbf{Statement I:} During zygotene, chromosomes start pairing together and this process of association is called synapsis.\\
\textbf{Statement II:} Electron micrographs of zygotene stage indicate that chromosome synapsis is accompanied by the formation of a complex structure called synaptonemal complex.\\
In the light of the above statements, choose the correct answer:
\begin{enumerate}[label=(\Alph*)]
    \item Both Statement I and Statement II are correct.
    \item Both Statement I and Statement II are incorrect.
    \item Statement I is correct but Statement II is incorrect.
    \item Statement I is incorrect but Statement II is correct.
\end{enumerate}

\item A diploid meiocyte has $2n = 24$ chromosomes and $2C$ amount of DNA in the $G_1$ phase. What will be the number of bivalents, tetrads, chromatids, and DNA content respectively during the pachytene stage of Meiosis I?
\begin{enumerate}[label=(\Alph*)]
    \item $12, 12, 48, 4C$
    \item $24, 12, 48, 4C$
    \item $12, 24, 24, 2C$
    \item $12, 12, 24, 4C$
\end{enumerate}

\item Which of the following features is clearly visible for the first time during the pachytene stage rather than the zygotene stage?
\begin{enumerate}[label=(\Alph*)]
    \item Homologous chromosomes starting to pair together
    \item Formation of the synaptonemal complex
    \item Four chromatids of each bivalent chromosome becoming distinct and clearly appearing as tetrads
    \item X-shaped structures called chiasmata at the sites of crossovers
\end{enumerate}

\item Regarding recombination nodules, which of the following statements is strictly accurate according to NCERT?
\begin{enumerate}[label=(\Alph*)]
    \item They appear during diplotene at the sites where crossing over has already completed.
    \item They appear during pachytene and represent the sites at which crossing over occurs between sister chromatids.
    \item They appear during pachytene and represent the sites at which crossing over occurs between non-sister chromatids of homologous chromosomes.
    \item They are formed during zygotene to initiate the formation of the synaptonemal complex.
\end{enumerate}

\item Crossing over is an enzyme-mediated process. Which of the following correctly describes its mechanism and outcome as stated in NCERT?
\begin{enumerate}[label=(\Alph*)]
    \item It is mediated by DNA polymerase and leads to recombination of genetic material on non-homologous chromosomes.
    \item It is mediated by recombinase and leads to recombination of genetic material on the two chromosomes.
    \item It occurs during diplotene and leaves chromosomes linked at the sites of crossovers until telophase I.
    \item It is completed by the end of zygotene, leaving the chromosomes attached at centromeres.
\end{enumerate}

\item Assertion (A): In oocytes of some vertebrates, diplotene can last for months or years.\\
Reason (R): The beginning of diplotene is recognised by the dissolution of the synaptonemal complex and the tendency of the recombined homologous chromosomes of the bivalents to separate from each other except at the sites of crossovers.
\begin{enumerate}[label=(\Alph*)]
    \item Both (A) and (R) are true and (R) is the correct explanation of (A).
    \item Both (A) and (R) are true but (R) is not the correct explanation of (A).
    \item (A) is true but (R) is false.
    \item (A) is false but (R) is true.
\end{enumerate}

\item Which stage of Meiosis I represents transition to metaphase?
\begin{enumerate}[label=(\Alph*)]
    \item Diplotene
    \item Diakinesis
    \item Pachytene
    \item Late Anaphase I
\end{enumerate}

\item All of the following events take place during diakinesis, \textbf{except}:
\begin{enumerate}[label=(\Alph*)]
    \item Terminalisation of chiasmata
    \item Chromosomes are fully condensed
    \item Meiotic spindle is assembled to prepare the homologous chromosomes for separation
    \item Dissolution of the synaptonemal complex begins
\end{enumerate}

\item By the end of diakinesis, which of the following cellular structures disappear?
\begin{enumerate}[label=(\Alph*)]
    \item Nucleolus only
    \item Nuclear envelope only
    \item Both nucleolus and nuclear envelope
    \item Centromeres and kinetochores
\end{enumerate}

\item How does Metaphase I of Meiosis differ fundamentally from Metaphase of Mitosis?
\begin{enumerate}[label=(\Alph*)]
    \item Chromosomes are maximally condensed in Mitosis but uncondensed in Meiosis I.
    \item In Metaphase I, bivalent chromosomes align on the equatorial plate, whereas in Mitosis, individual duplicated chromosomes align on the equatorial plate.
    \item Microtubules from opposite poles attach to the same homologous chromosome in Metaphase I.
    \item Centromeres split during Metaphase I but remain intact during Metaphase of Mitosis.
\end{enumerate}

\item During Metaphase I, the microtubules from the opposite poles of the spindle attach to:
\begin{enumerate}[label=(\Alph*)]
    \item The kinetochores of sister chromatids of the same chromosome.
    \item The kinetochore of homologous chromosomes.
    \item The chiasmata of the bivalent chromosomes.
    \item The telomeres of non-sister chromatids.
\end{enumerate}

\item Read the following statements about Anaphase I and select the correct option:
\begin{enumerate}[label=(\roman*)]
    \item The homologous chromosomes separate.
    \item Sister chromatids separate from each other.
    \item Sister chromatids remain associated at their centromeres.
    \item Splitting of centromere occurs simultaneously with chromosome movement.
\end{enumerate}
\begin{enumerate}[label=(\Alph*)]
    \item (i) and (iii) are correct
    \item (ii) and (iv) are correct
    \item (i) and (iv) are correct
    \item (ii) and (iii) are correct
\end{enumerate}

\item Reduction in chromosome number from diploid ($2n$) to haploid ($n$) actually occurs during which specific stage of Meiosis?
\begin{enumerate}[label=(\Alph*)]
    \item Metaphase I
    \item Anaphase I
    \item Anaphase II
    \item Telophase II
\end{enumerate}

\item During Telophase I of Meiosis:
\begin{enumerate}[label=(\Alph*)]
    \item The nuclear membrane and nucleolus reappear, cytokinesis follows, and this is called as dyad of cells.
    \item Chromosomes undergo complete decondensation to reach the extremely extended state of the interphase nucleus.
    \item Four haploid daughter cells are formed immediately after karyokinesis.
    \item DNA replication takes place inside the newly formed dyad nuclei.
\end{enumerate}

\item Consider the following statement from NCERT: "Although in many cases the chromosomes do undergo some dispersion, they do not reach the extremely extended state of the interphase nucleus." This statement refers to:
\begin{enumerate}[label=(\Alph*)]
    \item Diakinesis
    \item Telophase I
    \item Telophase II
    \item Prophase II
\end{enumerate}

\item The stage between the two meiotic divisions is called:
\begin{enumerate}[label=(\Alph*)]
    \item Interphase, which is long-lived and involves DNA replication.
    \item Interkinesis, which is generally short-lived and involves no DNA replication.
    \item Interkinesis, which is followed by Prophase I.
    \item $G_0$ phase, where cells remain metabolically active without dividing.
\end{enumerate}

\item A cell in $G_1$ phase has $16$ chromosomes and $10\text{ pg}$ of DNA. After Telophase I and cytokinesis, each cell of the dyad will have:
\begin{enumerate}[label=(\Alph*)]
    \item $16$ chromosomes and $10\text{ pg}$ of DNA
    \item $8$ chromosomes and $10\text{ pg}$ of DNA
    \item $8$ chromosomes and $5\text{ pg}$ of DNA
    \item $16$ chromosomes and $20\text{ pg}$ of DNA
\end{enumerate}

\item Which of the following events occurs during leptotene?
\begin{enumerate}[label=(\Alph*)]
    \item Compaction of chromosomes continues throughout leptotene.
    \item Pairing of homologous chromosomes begins.
    \item Appearance of recombination nodules.
    \item Formation of X-shaped chiasmata.
\end{enumerate}

\item Identify the incorrect match between the stage of Prophase I and its key event:
\begin{enumerate}[label=(\Alph*)]
    \item Zygotene — Synapsis
    \item Pachytene — Crossing over
    \item Diplotene — Dissolution of synaptonemal complex
    \item Diakinesis — Assembly of synaptonemal complex
\end{enumerate}

\item Why are the paired chromosomes in zygotene called either a bivalent or a tetrad?
\begin{enumerate}[label=(\Alph*)]
    \item Bivalent because there are two homologous chromosomes; tetrad because each chromosome has two chromatids, making four chromatids in total.
    \item Bivalent because there are two centromeres; tetrad because there are four centromeres.
    \item Bivalent because crossing over happens twice; tetrad because four daughter cells will be formed.
    \item Bivalent because two synaptonemal complexes form; tetrad because four recombination nodules appear.
\end{enumerate}

\item Which of the following statements is \textbf{not} true about crossing over?
\begin{enumerate}[label=(\Alph*)]
    \item It is the exchange of genetic material between two homologous chromosomes.
    \item It occurs between non-sister chromatids of homologous chromosomes.
    \item It is completed by the end of diplotene.
    \item It involves an enzyme called recombinase.
\end{enumerate}

\item At which stage are X-shaped structures called chiasmata visible?
\begin{enumerate}[label=(\Alph*)]
    \item Leptotene
    \item Zygotene
    \item Pachytene
    \item Diplotene
\end{enumerate}

\item If a diploid organism has $2n = 36$, how many chiasmata at minimum would be observed during diplotene if every bivalent undergoes at least one crossover?
\begin{enumerate}[label=(\Alph*)]
    \item $36$
    \item $18$
    \item $72$
    \item $9$
\end{enumerate}

\item In Meiosis I, the chromosomes that pair during zygotene are:
\begin{enumerate}[label=(\Alph*)]
    \item Non-homologous chromosomes of identical length
    \item Homologous chromosomes
    \item Sister chromatids of a single chromosome
    \item Sex chromosomes only
\end{enumerate}

\item Assertion (A): During Anaphase I, the centromeres do not split.\\
Reason (R): In Anaphase I, homologous chromosomes separate while sister chromatids remain associated at their centromeres.
\begin{enumerate}[label=(\Alph*)]
    \item Both (A) and (R) are true and (R) is the correct explanation of (A).
    \item Both (A) and (R) are true but (R) is not the correct explanation of (A).
    \item (A) is true but (R) is false.
    \item (A) is false but (R) is true.
\end{enumerate}

\item Which of the following is a distinguishing feature of Prophase I compared to Prophase II?
\begin{enumerate}[label=(\Alph*)]
    \item Disappearance of nuclear membrane
    \item Compaction of chromosomes
    \item Pairing of homologous chromosomes and recombination between non-sister chromatids
    \item Formation of spindle fibres
\end{enumerate}

\item During which phase of Meiosis I is the meiotic spindle assembled to prepare the homologous chromosomes for separation?
\begin{enumerate}[label=(\Alph*)]
    \item Pachytene
    \item Diplotene
    \item Diakinesis
    \item Metaphase I
\end{enumerate}

\item Given below are two statements:
\textbf{Statement I:} Meiosis I is initiated after the parental chromosomes have replicated to produce identical sister chromatids at the S phase.\\
\textbf{Statement II:} Meiosis involves two sequential cycles of nuclear and cell division called meiosis I and meiosis II, as well as two cycles of DNA replication.\\
Choose the most appropriate answer:
\begin{enumerate}[label=(\Alph*)]
    \item Both Statement I and Statement II are correct.
    \item Both Statement I and Statement II are incorrect.
    \item Statement I is correct but Statement II is incorrect.
    \item Statement I is incorrect but Statement II is correct.
\end{enumerate}

\item How many spindle fibres (microtubules from poles) attach to a single bivalent during Metaphase I according to NCERT?
\begin{enumerate}[label=(\Alph*)]
    \item One microtubule from one pole only
    \item Microtubules from opposite poles attach to the kinetochore of homologous chromosomes
    \item Four microtubules attach to four separate centromeres
    \item Microtubules attach directly to the chiasmata
\end{enumerate}

\item Which of the following correctly arranges the events of Prophase I in chronological order?
\begin{enumerate}[label=(\roman*)]
    \item Appearance of recombination nodules
    \item Dissolution of synaptonemal complex
    \item Terminalisation of chiasmata
    \item Compaction of chromosomes makes them gradually visible under light microscope
    \item Formation of synaptonemal complex
\end{enumerate}
\begin{enumerate}[label=(\Alph*)]
    \item (iv) $\rightarrow$ (v) $\rightarrow$ (i) $\rightarrow$ (ii) $\rightarrow$ (iii)
    \item (iv) $\rightarrow$ (i) $\rightarrow$ (v) $\rightarrow$ (ii) $\rightarrow$ (iii)
    \item (v) $\rightarrow$ (iv) $\rightarrow$ (i) $\rightarrow$ (iii) $\rightarrow$ (ii)
    \item (iv) $\rightarrow$ (v) $\rightarrow$ (ii) $\rightarrow$ (i) $\rightarrow$ (iii)
\end{enumerate}

\item In a hypothetical mutant meiocyte where the enzyme \textit{recombinase} is non-functional, which of the following events will still occur normally?
\begin{enumerate}[label=(\Alph*)]
    \item Exchange of genetic material between non-sister chromatids
    \item Formation of synaptonemal complex during zygotene
    \item Formation of chiasmata due to crossovers
    \item Genetic recombination on homologous chromosomes
\end{enumerate}

\item Which sub-stage of Prophase I is characterised by the disappearance of the nucleolus and breakdown of the nuclear envelope?
\begin{enumerate}[label=(\Alph*)]
    \item Leptotene
    \item Zygotene
    \item Diplotene
    \item Diakinesis
\end{enumerate}

\item If a meiocyte has $20$ bivalents at Metaphase I, how many chromosomes and chromatids respectively will be present at each pole at the end of Anaphase I?
\begin{enumerate}[label=(\Alph*)]
    \item $20$ chromosomes and $40$ chromatids
    \item $10$ chromosomes and $20$ chromatids
    \item $20$ chromosomes and $20$ chromatids
    \item $40$ chromosomes and $40$ chromatids
\end{enumerate}

\item Which of the following statements about the synaptonemal complex is true as per NCERT?
\begin{enumerate}[label=(\Alph*)]
    \item It is observed using a light microscope during leptotene.
    \item It is observed in electron micrographs of the zygotene stage.
    \item It is formed during pachytene to hold sister chromatids together.
    \item It persists throughout Metaphase I until Anaphase I begins.
\end{enumerate}

\item By the end of which stage is crossing over completed, leaving the chromosomes linked at the sites of crossing over?
\begin{enumerate}[label=(\Alph*)]
    \item Zygotene
    \item Pachytene
    \item Diplotene
    \item Diakinesis
\end{enumerate}

\item During diplotene, the recombined homologous chromosomes of the bivalents have a tendency to separate from each other except at:
\begin{enumerate}[label=(\Alph*)]
    \item The centromeres
    \item The telomeres
    \item The sites of crossovers
    \item The kinetochores
\end{enumerate}

\item What is the ploidy level of each cell in the "dyad of cells" formed after Telophase I and cytokinesis?
\begin{enumerate}[label=(\Alph*)]
    \item Diploid ($2n$) with $2C$ DNA
    \item Haploid ($n$) with $2C$ DNA
    \item Haploid ($n$) with $1C$ DNA
    \item Diploid ($2n$) with $4C$ DNA
\end{enumerate}

\item Compare Metaphase I of Meiosis and Metaphase of Mitosis. Which statement is accurate?
\begin{enumerate}[label=(\Alph*)]
    \item Both involve alignment of bivalents at the equator.
    \item In Metaphase I, two centromeres of each bivalent face opposite poles, whereas in Mitotic Metaphase, sister chromatids of each chromosome face opposite poles.
    \item Both involve splitting of centromeres at the end of metaphase.
    \item Spindle fibres are absent in Metaphase I.
\end{enumerate}

\item Which of the following is \textbf{not} a characteristic of interkinesis?
\begin{enumerate}[label=(\Alph*)]
    \item It is the stage between the two meiotic divisions.
    \item It is generally short-lived.
    \item There is no replication of DNA during this stage.
    \item It is followed by Prophase I.
\end{enumerate}

\item Select the correct combination of statements regarding Pachytene:
\begin{enumerate}[label=(\roman*)]
    \item Bivalent chromosomes clearly appear as tetrads.
    \item This stage is characterised by the appearance of recombination nodules.
    \item Crossing over occurs between sister chromatids of homologous chromosomes.
    \item Recombination between homologous chromosomes is completed by the end of pachytene.
\end{enumerate}
\begin{enumerate}[label=(\Alph*)]
    \item (i), (ii), and (iv) only
    \item (i), (ii), and (iii) only
    \item (ii), (iii), and (iv) only
    \item (i) and (iii) only
\end{enumerate}

\item During Anaphase I, the separation of homologous chromosomes is termed disjunction. What happens to the DNA content per chromosome during Anaphase I compared to Metaphase I?
\begin{enumerate}[label=(\Alph*)]
    \item It is halved because chromatids separate.
    \item It remains unchanged because each chromosome still consists of two chromatids attached at a single centromere.
    \item It doubles due to DNA synthesis.
    \item It becomes zero until Telophase I.
\end{enumerate}

\item Which of the following pairs of stages and events is completely accurate according to NCERT?
\begin{enumerate}[label=(\Alph*)]
    \item Leptotene — Bivalent formation
    \item Diplotene — Terminalisation of chiasmata
    \item Diakinesis — Represents transition to metaphase
    \item Telophase I — Chromosomes reach extremely extended state of interphase
\end{enumerate}

\end{enumerate}

\newpage
\section*{Part II: Answer Key}

\begin{multicols}{5}
\begin{enumerate}[label=\textbf{\arabic*.}]
    \item (A)
    \item (D)
    \item (A)
    \item (A)
    \item (A)
    \item (C)
    \item (C)
    \item (B)
    \item (B)
    \item (B)
    \item (D)
    \item (C)
    \item (B)
    \item (B)
    \item (A)
    \item (B)
    \item (A)
    \item (B)
    \item (B)
    \item (B)
    \item (A)
    \item (D)
    \item (A)
    \item (C)
    \item (D)
    \item (B)
    \item (B)
    \item (A)
    \item (C)
    \item (C)
    \item (C)
    \item (B)
    \item (A)
    \item (B)
    \item (D)
    \item (A)
    \item (B)
    \item (B)
    \item (C)
    \item (B)
    \item (B)
    \item (D)
    \item (A)
    \item (B)
    \item (C)
\end{enumerate}
\end{multicols}

\newpage
\section*{Part III: Detailed Solutions (NCERT Reference)}

\begin{enumerate}[label=\textbf{\arabic*.}]

\item \textbf{Answer: (A)}\\
\textbf{Explanation:} In Meiosis I, homologous chromosomes pair during zygotene (\textit{synapsis}), exchange segments during pachytene (\textit{crossing over}), undergo \textit{terminalisation of chiasmata} during diakinesis, and finally separate during Anaphase I (\textit{disjunction of homologs}). Centromeres do not split in Meiosis I.

\item \textbf{Answer: (D)}\\
\textbf{Explanation:} According to NCERT (Class 11, Chapter 10, Section 10.4.1), during leptotene the chromosomes become gradually visible under the light microscope, and the compaction of chromosomes \textit{continues throughout leptotene}. It does not stop at the end of leptotene.

\item \textbf{Answer: (A)}\\
\textbf{Explanation:}
\begin{itemize}
    \item \textbf{Leptotene:} Compaction of chromosomes continues throughout leptotene (iii).
    \item \textbf{Zygotene:} Chromosomes start pairing (synapsis) to form a bivalent or tetrad (iv).
    \item \textbf{Pachytene:} Crossing over mediated by the enzyme recombinase occurs (i).
    \item \textbf{Diplotene:} Recognised by the dissolution of the synaptonemal complex (ii).
\end{itemize}

\item \textbf{Answer: (A)}\\
\textbf{Explanation:} Both statements are verbatim facts from NCERT regarding the zygotene stage of Prophase I.

\item \textbf{Answer: (A)}\\
\textbf{Explanation:} Given $2n = 24$ and $2C$ DNA in $G_1$:
\begin{itemize}
    \item After S phase, DNA doubles to $4C$, while chromosome number remains $24$ (each having $2$ chromatids, so $48$ chromatids total).
    \item During zygotene/pachytene, $24$ homologous chromosomes pair into $12$ bivalents (each bivalent is also called a tetrad, hence $12$ tetrads).
    \item Total chromatids $= 12 \times 4 = 48$; DNA content $= 4C$.
\end{itemize}

\item \textbf{Answer: (C)}\\
\textbf{Explanation:} NCERT explicitly notes that while the complex formed by a pair of synapsed homologous chromosomes in zygotene is called a bivalent or a tetrad, \textit{"these are more clearly visible at the next stage [pachytene]"} where the four chromatids of each bivalent chromosome become distinct and clearly appear as tetrads.

\item \textbf{Answer: (C)}\\
\textbf{Explanation:} Pachytene is characterised by the appearance of recombination nodules, which are the sites at which crossing over occurs between \textit{non-sister chromatids} of the homologous chromosomes.

\item \textbf{Answer: (B)}\\
\textbf{Explanation:} NCERT states: \textit{"Crossing over is also an enzyme-mediated process and the enzyme involved is called recombinase. Crossing over leads to recombination of genetic material on the two chromosomes."}

\item \textbf{Answer: (B)}\\
\textbf{Explanation:} Both statements are true NCERT statements about diplotene. However, the dissolution of the synaptonemal complex describes how diplotene begins, rather than explaining the physiological arrest of vertebrate oocytes (which serves for yolk synthesis/growth, not detailed in NCERT as caused by SC dissolution).

\item \textbf{Answer: (B)}\\
\textbf{Explanation:} NCERT explicitly states that \textit{diakinesis} represents transition to metaphase.

\item \textbf{Answer: (D)}\\
\textbf{Explanation:} Dissolution of the synaptonemal complex marks the beginning of \textit{diplotene}, whereas terminalisation of chiasmata, full condensation of chromosomes, and assembly of the meiotic spindle occur during \textit{diakinesis}.

\item \textbf{Answer: (C)}\\
\textbf{Explanation:} By the end of diakinesis, the nucleolus disappears and the nuclear envelope also breaks down.

\item \textbf{Answer: (B)}\\
\textbf{Explanation:} In Metaphase I, the bivalent chromosomes align on the equatorial plate, whereas in mitotic metaphase, individual duplicated chromosomes align on the metaphase plate.

\item \textbf{Answer: (B)}\\
\textbf{Explanation:} NCERT states: \textit{"The microtubules from the opposite poles of the spindle attach to the kinetochore of homologous chromosomes."}

\item \textbf{Answer: (A)}\\
\textbf{Explanation:} During Anaphase I, the homologous chromosomes separate, while sister chromatids remain associated at their centromeres. Splitting of centromeres occurs in Anaphase II, not Anaphase I.

\item \textbf{Answer: (B)}\\
\textbf{Explanation:} Reduction in chromosome number from $2n$ to $n$ happens during Anaphase I when homologous chromosomes segregate to opposite poles.

\item \textbf{Answer: (A)}\\
\textbf{Explanation:} In Telophase I, the nuclear membrane and nucleolus reappear, cytokinesis follows, and this stage is called a \textit{dyad of cells}.

\item \textbf{Answer: (B)}\\
\textbf{Explanation:} NCERT states this specifically for Telophase I, noting that chromosomes undergo some dispersion without reaching the extremely extended state of an interphase nucleus.

\item \textbf{Answer: (B)}\\
\textbf{Explanation:} The stage between the two meiotic divisions is called \textit{interkinesis} and is generally short-lived. There is no replication of DNA during interkinesis, and it is followed by Prophase II.

\item \textbf{Answer: (B)}\\
\textbf{Explanation:} In $G_1$: $2n = 16$, DNA $= 10\text{ pg}$. After S phase: $2n = 16$, DNA $= 20\text{ pg}$. After Meiosis I (Telophase I + cytokinesis), homologous chromosomes separate into two cells without centromere splitting, so each cell receives $n = 8$ chromosomes and half the duplicated DNA ($10\text{ pg}$).

\item \textbf{Answer: (A)}\\
\textbf{Explanation:} Compaction of chromosomes continues throughout leptotene, making them gradually visible under the light microscope.

\item \textbf{Answer: (D)}\\
\textbf{Explanation:} Diakinesis involves assembly of the \textit{meiotic spindle}, not the synaptonemal complex (which forms in zygotene and dissolves in diplotene).

\item \textbf{Answer: (A)}\\
\textbf{Explanation:} A pair of synapsed homologous chromosomes consists of two chromosomes (bivalent) and four chromatids in total (tetrad).

\item \textbf{Answer: (C)}\\
\textbf{Explanation:} Recombination between homologous chromosomes is completed by the end of \textit{pachytene}, leaving the chromosomes linked at the sites of crossing over.

\item \textbf{Answer: (D)}\\
\textbf{Explanation:} During diplotene, the separation of homologous chromosomes except at the sites of crossovers forms X-shaped structures called chiasmata.

\item \textbf{Answer: (B)}\\
\textbf{Explanation:} With $2n = 36$ chromosomes, there are $18$ homologous pairs (bivalents). If each bivalent undergoes at least one crossover, there will be a minimum of $18$ chiasmata visible in diplotene.

\item \textbf{Answer: (B)}\\
\textbf{Explanation:} During zygotene, homologous chromosomes start pairing together (synapsis).

\item \textbf{Answer: (A)}\\
\textbf{Explanation:} Because homologous chromosomes separate as intact dyads (sister chromatids remain attached at their centromeres), the centromeres do not split during Anaphase I.

\item \textbf{Answer: (C)}\\
\textbf{Explanation:} Pairing of homologous chromosomes (synapsis) and recombination between non-sister chromatids occur only in Prophase I, not in Prophase II (which resembles normal mitotic prophase).

\item \textbf{Answer: (C)}\\
\textbf{Explanation:} During diakinesis, the chromosomes are fully condensed and the meiotic spindle is assembled to prepare the homologous chromosomes for separation.

\item \textbf{Answer: (C)}\\
\textbf{Explanation:} Statement I is correct. Statement II is incorrect because meiosis involves two sequential cycles of nuclear and cell division, but \textit{only a single cycle of DNA replication}.

\item \textbf{Answer: (B)}\\
\textbf{Explanation:} NCERT specifies that during Metaphase I, the microtubules from the opposite poles of the spindle attach to the kinetochore of homologous chromosomes.

\item \textbf{Answer: (A)}\\
\textbf{Explanation:} The correct order is: Leptotene (iv) $\rightarrow$ Zygotene (v) $\rightarrow$ Pachytene (i) $\rightarrow$ Diplotene (ii) $\rightarrow$ Diakinesis (iii).

\item \textbf{Answer: (B)}\\
\textbf{Explanation:} Formation of the synaptonemal complex occurs during zygotene prior to recombinase-mediated crossing over in pachytene. Thus, synapsis and synaptonemal complex assembly would proceed normally.

\item \textbf{Answer: (D)}\\
\textbf{Explanation:} By the end of diakinesis, the nucleolus disappears and the nuclear envelope breaks down.

\item \textbf{Answer: (A)}\\
\textbf{Explanation:} $20$ bivalents equal $40$ chromosomes ($2n = 40$). During Anaphase I, homologous chromosomes separate, so $20$ chromosomes move to each pole. Since centromeres do not split, each chromosome still has $2$ chromatids, giving $20 \times 2 = 40$ chromatids at each pole.

\item \textbf{Answer: (B)}\\
\textbf{Explanation:} NCERT states: \textit{"Electron micrographs of this stage [zygotene] indicate that chromosome synapsis is accompanied by the formation of complex structure called synaptonemal complex."}

\item \textbf{Answer: (B)}\\
\textbf{Explanation:} Recombination between homologous chromosomes is completed by the end of pachytene, leaving the chromosomes linked at the sites of crossing over.

\item \textbf{Answer: (C)}\\
\textbf{Explanation:} In diplotene, recombined homologous chromosomes tend to separate from each other except at the sites of crossovers, forming X-shaped chiasmata.

\item \textbf{Answer: (B)}\\
\textbf{Explanation:} After Telophase I and cytokinesis, the chromosome number is reduced to haploid ($n$), while each chromosome still possesses two sister chromatids, keeping the DNA content at $2C$.

\item \textbf{Answer: (B)}\\
\textbf{Explanation:} In Metaphase I, homologous chromosomes (two centromeres per bivalent) face opposite poles. In mitotic metaphase, sister chromatids of each single chromosome attach to spindle fibres from opposite poles.

\item \textbf{Answer: (D)}\\
\textbf{Explanation:} Interkinesis is followed by \textit{Prophase II}, not Prophase I.

\item \textbf{Answer: (A)}\\
\textbf{Explanation:} Statement (iii) is incorrect because crossing over occurs between \textit{non-sister} chromatids, not sister chromatids. Statements (i), (ii), and (iv) are completely accurate.

\item \textbf{Answer: (B)}\\
\textbf{Explanation:} During Anaphase I, homologous chromosomes separate without centromere division, so each chromosome retains both sister chromatids and its DNA content remains unchanged.

\item \textbf{Answer: (C)}\\
\textbf{Explanation:} Diakinesis represents transition to metaphase. Bivalent formation is zygotene, terminalisation is diakinesis, and Telophase I chromosomes do not reach the extremely extended interphase state.

\end{enumerate}

\end{document}